\documentclass[12pt]{article}
\usepackage{amsmath, amssymb}
\usepackage[margin=1in]{geometry}
\setlength{\parskip}{6pt}
\title{QUBO Formulation: FPGA Placement Problem}
\date{}
\begin{document}
\maketitle

\subsection*{FPGA Placement Problem}

\paragraph{Problem Statement.}
Given a set of $N$ logic blocks $\mathcal{B} = \{0, 1, \dots, N-1\}$ and a set of $M$ discrete FPGA sites $\mathcal{S} = \{0, 1, \dots, M-1\}$, the goal is to assign each block to a unique site such that the total weighted communication distance between connected blocks is minimized. Each site can accommodate at most one block. A set of connections $\mathcal{C} \subseteq \mathcal{B} \times \mathcal{B}$ defines pairs of blocks that communicate, with associated weights $w_{i,k}$ representing the importance or frequency of communication between block $i$ and block $k$. The distance between any two sites $j, l \in \mathcal{S}$ is given by a metric $d_{j,l} \ge 0$. The objective is to find a placement that minimizes the sum of weighted distances between all connected block pairs based on their assigned sites.

\paragraph{Decision Variables.}
We introduce binary variables $x_{i,j} \in \{0, 1\}$ for each block $i \in \mathcal{B}$ and site $j \in \mathcal{S}$. The variable $x_{i,j}$ takes the value $1$ if and only if block $i$ is assigned to site $j$, and $0$ otherwise. The total number of variables is $n = N \times M$.

\paragraph{QUBO Derivation.}
We construct a Quadratic Unconstrained Binary Optimization (QUBO) Hamiltonian $H(\mathbf{x})$ by encoding the objective function and constraints into a single energy function. The QUBO form is $H(\mathbf{x}) = \sum_{u \le v} Q_{u,v} x_u x_v + c$, where $x_u$ denotes the variable indexed by $u = i M + j$.

\textbf{Step 1: Objective Function.}
The objective is to minimize the total weighted wirelength:
\begin{equation}
    E_{\text{obj}} = \sum_{(i,k) \in \mathcal{C}} \sum_{j \in \mathcal{S}} \sum_{l \in \mathcal{S}} w_{i,k} \, d_{j,l} \, x_{i,j} \, x_{k,l}.
\end{equation}
Since QUBO minimizes energy, we set $H_{\text{obj}} = E_{\text{obj}}$.

\textbf{Step 2: Constraints.}
The problem imposes two constraints:
\begin{enumerate}
    \item \textit{Block Assignment:} Each block must be assigned to exactly one site:
    \begin{equation}
        \sum_{j \in \mathcal{S}} x_{i,j} = 1, \quad \forall i \in \mathcal{B}.
    \end{equation}
    \item \textit{Site Capacity:} Each site can hold at most one block:
    \begin{equation}
        \sum_{i \in \mathcal{B}} x_{i,j} \le 1, \quad \forall j \in \mathcal{S}.
    \end{equation}
\end{enumerate}

\textbf{Step 3: Penalty Terms.}
We convert constraints to penalties using a weight $\lambda > 0$.

\textit{Block Assignment Penalty:}
We square the constraint violation:
\begin{align}
    P_{\text{block}} &= \lambda \sum_{i \in \mathcal{B}} \left( \sum_{j \in \mathcal{S}} x_{i,j} - 1 \right)^2 \\
    &= \lambda \sum_{i \in \mathcal{B}} \left( \sum_{j \in \mathcal{S}} x_{i,j}^2 + \sum_{j \neq l \in \mathcal{S}} x_{i,j} x_{i,l} - 2 \sum_{j \in \mathcal{S}} x_{i,j} + 1 \right).
\end{align}
Using the idempotent property $x_{i,j}^2 = x_{i,j}$ and noting that $\sum_{j \neq l} x_{i,j} x_{i,l} = 2 \sum_{j < l \in \mathcal{S}} x_{i,j} x_{i,l}$, we expand:
\begin{align}
    P_{\text{block}} &= \lambda \sum_{i \in \mathcal{B}} \left( \sum_{j \in \mathcal{S}} x_{i,j} + 2 \sum_{j < l \in \mathcal{S}} x_{i,j} x_{i,l} - 2 \sum_{j \in \mathcal{S}} x_{i,j} + 1 \right) \\
    &= \lambda \sum_{i \in \mathcal{B}} \left( 1 - \sum_{j \in \mathcal{S}} x_{i,j} + 2 \sum_{j < l \in \mathcal{S}} x_{i,j} x_{i,l} \right).
\end{align}
This yields linear terms $-\lambda x_{i,j}$, quadratic terms $2\lambda x_{i,j} x_{i,l}$ for $j \neq l$, and a constant offset contribution $\lambda$.

\textit{Site Capacity Penalty:}
The inequality $\sum_{i} x_{i,j} \le 1$ is enforced by penalizing pairs of blocks assigned to the same site:
\begin{equation}
    P_{\text{site}} = \lambda \sum_{j \in \mathcal{S}} \sum_{i < i' \in \mathcal{B}} x_{i,j} x_{i',j}.
\end{equation}
This term is zero if at most one block is assigned to site $j$, and positive otherwise.

\textbf{Step 4: Combined Hamiltonian.}
The total QUBO Hamiltonian is:
\begin{equation}
    H(\mathbf{x}) = H_{\text{obj}} + P_{\text{block}} + P_{\text{site}}.
\end{equation}

\textbf{Step 5: Q Matrix Entries and Offset.}
Reading off the coefficients from the expanded form, the QUBO parameters are:
\begin{align}
    Q_{(i,j),(k,l)} &= w_{i,k} \, d_{j,l}, & \forall (i,k) \in \mathcal{C}, \, j,l \in \mathcal{S}, \\
    Q_{(i,j),(i,j)} &= -\lambda, & \forall i \in \mathcal{B}, \, j \in \mathcal{S}, \\
    Q_{(i,j),(i,l)} &= 2\lambda, & \forall i \in \mathcal{B}, \, j,l \in \mathcal{S}, \, j \neq l, \\
    Q_{(i,j),(i',j)} &= \lambda, & \forall i,i' \in \mathcal{B}, \, j \in \mathcal{S}, \, i \neq i', \\
    c &= \lambda N.
\end{align}
All other entries $Q_{u,v}$ are zero.

\paragraph{Penalty Weight Justification.}
The penalty weight $\lambda$ must be chosen sufficiently large to ensure that any constraint violation increases the energy more than any possible reduction in the objective function. A sufficient condition is:
\begin{equation}
    \lambda > \sum_{(i,k) \in \mathcal{C}} |w_{i,k}|.
\end{equation}
This guarantees that the minimum penalty cost for violating a constraint exceeds the maximum possible gain achievable by reducing the objective, thereby ensuring that the global minimum of $H(\mathbf{x})$ corresponds to a feasible solution.

\paragraph{Correctness Argument.}
Minimizing $H(\mathbf{x})$ over $\{0,1\}^n$ recovers the optimal placement because: (i) for any feasible assignment, the penalty terms vanish, so $H(\mathbf{x}) = E_{\text{obj}}(\mathbf{x})$; (ii) for any infeasible assignment, the penalty terms are strictly positive. Given $\lambda > \sum |w_{i,k}|$, the energy of any infeasible solution exceeds that of the best feasible solution, as the penalty increase dominates any potential decrease in $E_{\text{obj}}$. Thus, the global minimum of $H(\mathbf{x})$ must be feasible and minimize $E_{\text{obj}}$, yielding the optimal placement.

\paragraph{Illustrative Example.}
Consider a small instance with $N=2$ blocks $\mathcal{B}=\{0,1\}$ and $M=2$ sites $\mathcal{S}=\{0,1\}$. There is a single connection $(0,1)$ with weight $w_{0,1}=1$. The distance metric is $d_{0,1}=d_{1,0}=1$ and $d_{0,0}=d_{1,1}=0$. We set $\lambda=10$.

The variables are $x_{0,0}, x_{0,1}, x_{1,0}, x_{1,1}$ mapped to indices $0,1,2,3$. The objective terms are $x_{0,0}x_{1,1}$ and $x_{0,1}x_{1,0}$, each with coefficient $1$.

The block assignment penalty expands to:
\begin{align}
    P_{\text{block}} &= 10 \left[ (x_{0,0}+x_{0,1}-1)^2 + (x_{1,0}+x_{1,1}-1)^2 \right] \\
    &= 20 - 10(x_{0,0}+x_{0,1}+x_{1,0}+x_{1,1}) + 20(x_{0,0}x_{0,1} + x_{1,0}x_{1,1}).
\end{align}
The site capacity penalty is:
\begin{equation}
    P_{\text{site}} = 10(x_{0,0}x_{1,0} + x_{0,1}x_{1,1}).
\end{equation}
The total Hamiltonian is:
\begin{align}
    H(\mathbf{x}) &= x_{0,0}x_{1,1} + x_{0,1}x_{1,0} + 20 - 10(x_{0,0}+x_{0,1}+x_{1,0}+x_{1,1}) \\
    &\quad + 20(x_{0,0}x_{0,1} + x_{1,0}x_{1,1}) + 10(x_{0,0}x_{1,0} + x_{0,1}x_{1,1}).
\end{align}
The constant offset is $c=20$. The optimal feasible solutions are $x_{0,0}=x_{1,1}=1$ (blocks at sites $0,1$) and $x_{0,1}=x_{1,0}=1$ (blocks at sites $1,0$), both yielding energy $H=1$. Any infeasible assignment incurs a penalty of at least $10$, resulting in energy $\ge 10$.

\end{document}